\documentclass[11pt,a4paper]{article}
\usepackage[T1]{fontenc}
\usepackage{lmodern}
\usepackage[margin=22mm,headheight=15pt]{geometry}
\usepackage{amsmath,amssymb,mathtools,microtype}
\usepackage{xcolor,enumitem,fancyhdr,booktabs,needspace}
\usepackage[hidelinks]{hyperref}
\definecolor{examblue}{HTML}{17365D}
\pagestyle{fancy}
\fancyhf{}
\lhead{\textsc{Machine Learning 1}}
\rhead{\textsc{Mock 2 --- A: Answer Key}}
\cfoot{\thepage}
\setlength{\parindent}{0pt}
\setlength{\parskip}{0.4em}
\setlist[itemize]{leftmargin=1.5em,itemsep=0.1em}
\newcommand{\R}{\mathbb R}
\newcommand{\E}{\mathbb E}
\DeclareMathOperator{\Cov}{Cov}
\DeclareMathOperator{\Var}{Var}
\DeclareMathOperator{\tr}{tr}
\newcommand{\solution}[3]{\Needspace{6\baselineskip}\subsection*{\color{examblue}#1\quad #2\hfill\normalsize[#3 p]}}
\newcommand{\selfcredit}[1]{\par\textit{Suggested credit: #1}\par}

\begin{document}
\begin{center}
{\Huge\bfseries\color{examblue}Mock Midterm 2 --- A}\par
\vspace{0.5em}
{\LARGE Worked solutions and marking guide}\par
\vspace{0.7em}
{\large 18 points: Exercise 1 (10), Exercise 2 (8)}
\end{center}
These are full solutions, not only final answers. The marking allocations are for
self-assessment, not an official course rubric. Equivalent derivations and the
opposite sign of a PCA eigenvector are valid. Give partial credit for a correct
method even if a later arithmetic error propagates.

\section*{Exercise 1: Two sensors, one class-dependent rate}
\solution{1a}{Joint likelihood and constrained estimation}{3}
First use the product rule $p(x,C_k)=\pi_kp(x\mid C_k)$, then conditional
independence of the two sensors. Independence of training observations gives
\[
p(X,T\mid\pi,\lambda)=\prod_{n=1}^N\prod_{k=1}^K
\left[\pi_k\frac{2^{m_n}}{m_n!}\lambda_k^{m_n+1}
e^{-\lambda_k(2+s_n)}\right]^{t_{nk}}.
\]
In particular, the two rate factors combine to $\lambda_k^{m_n+1}$.
Writing $B=\sum_n[m_n\log2-\log(m_n!)]$, which is parameter-independent,
\[
\ell=B+\sum_k\left[N_k\log\pi_k+(M_k+N_k)\log\lambda_k
-(2N_k+S_k)\lambda_k\right].
\]
For the rates,
\[
\frac{\partial\ell}{\partial\lambda_k}
=\frac{M_k+N_k}{\lambda_k}-(2N_k+S_k)=0
\quad\Longrightarrow\quad
\boxed{\widehat\lambda_k=\frac{M_k+N_k}{2N_k+S_k}}.
\]
The second derivative is $-(M_k+N_k)/\lambda_k^2<0$. The log-likelihood tends
to $-\infty$ at both ends of the positive rate domain, so this is its unique
global rate maximum. The positive count $N_k$ ensures both numerator and
denominator are positive, even when every count or waiting time in a class is zero.

For the priors use
\[
\mathcal L=\sum_kN_k\log\pi_k+\eta\left(\sum_k\pi_k-1\right).
\]
The stationary conditions give $N_k/\pi_k+\eta=0$, hence
$\pi_k=-N_k/\eta$. Summing gives $\eta=-N$ and therefore
\[
\boxed{\widehat\pi_k=N_k/N}.
\]
The prior part is strictly concave on the simplex interior and tends to
$-\infty$ if any $\pi_k\to0$ with $N_k>0$.
\selfcredit{1 p for the likelihood and log-likelihood; 1 p for the rate estimate and maximum check; 1 p for the constrained prior estimate.}

\solution{1b}{Posterior scores and class region}{2}
The fitted parameters are
\[
\boxed{(\widehat\pi_1,\widehat\pi_2,\widehat\pi_3)
=(\tfrac12,\tfrac14,\tfrac14),\qquad
(\widehat\lambda_1,\widehat\lambda_2,\widehat\lambda_3)=(1,2,\tfrac12)}.
\]
Bayes' rule gives
\[
p(C_k\mid m,s)=
\frac{\pi_k\lambda_k^{m+1}e^{-\lambda_k(2+s)}}
{\sum_j\pi_j\lambda_j^{m+1}e^{-\lambda_j(2+s)}}.
\]
The common factor $2^m/m!$ cancels. An affine score is
\[
g_k(m,s)=\log\pi_k+\log\lambda_k-2\lambda_k
+m\log\lambda_k-s\lambda_k,
\qquad p(C_k\mid m,s)=\operatorname{softmax}(g)_k.
\]
Thus $C_1$ is at least as probable as its two competitors precisely when
\begin{align*}
g_1-g_2&=2+s-m\log2\geq0,\\
g_1-g_3&=(m+2)\log2-1-\tfrac12s\geq0.
\end{align*}
Equivalently,
\[
\boxed{m\log2\leq2+s,\qquad 1+\tfrac12s\leq(m+2)\log2}.
\]
With $m,s\geq0$ treated as real coordinates, this region is an intersection
of half-spaces and is convex. The actual count-domain subset is not a convex
subset of $\R^2$ merely because it is restricted to integer $m$; the question
explicitly asks about the continuous relaxation. Ties may be assigned by any
specified tie rule, whereas these non-strict inequalities describe posterior dominance.
\selfcredit{0.5 p for parameters; 0.5 p for scores/posterior; 0.75 p for inequalities; 0.25 p for convexity with the domain qualification.}

\solution{1c}{A missing sensor changes the decision}{1}
For the complete input $(0,\tfrac12)$,
\[
g_1-g_2=\tfrac52>0,\qquad
g_1-g_3=2\log2-\tfrac54>0.
\]
Consequently the predicted class is $\boxed{C_1}$.

If the time is unobserved, integrate out the time density:
\[
p(m\mid C_k)=\int_0^\infty p(m,s\mid C_k)\,ds
=\frac{(2\lambda_k)^me^{-2\lambda_k}}{m!},
\]
since the exponential density integrates to one. Missingness itself gives no
additional class evidence under the stipulated independent failure mechanism.
The count-only posterior is
\[
\boxed{p(C_k\mid m)=
\frac{\pi_k\lambda_k^m e^{-2\lambda_k}}
{\sum_j\pi_j\lambda_j^m e^{-2\lambda_j}}}.
\]
At $m=0$ the unnormalized probabilities are
\[
\tfrac12e^{-2},\qquad\tfrac14e^{-4},\qquad\tfrac14e^{-1}.
\]
The last is largest: its ratio to the first is $e/2>1$. The new prediction is
$\boxed{C_3}$. Setting $s=0$ retains the factor $p(s=0\mid C_k)=\lambda_k$
as a density value and gives the complete-input prediction $C_1$, not the
count-only prediction. A missing observation is not an observed zero.
\selfcredit{0.5 p for marginalization and the count-only posterior; 0.5 p for both decisions and the zero-imputation distinction.}

\solution{1d}{Conditional independence is not marginal independence}{1}
Within a class,
\[
\E[MS\mid C_k]=\E[M\mid C_k]\E[S\mid C_k]=2.
\]
Thus $\E[MS]=2$, while
\[
\E[M]=2\sum_k\pi_k\lambda_k,
\qquad \E[S]=\sum_k\frac{\pi_k}{\lambda_k}.
\]
In general this gives
\[
\boxed{\Cov(M,S)=2-2\left(\sum_k\pi_k\lambda_k\right)
\left(\sum_k\frac{\pi_k}{\lambda_k}\right)}.
\]
For these fitted parameters,
\[
\sum_k\pi_k\lambda_k=\tfrac98,\qquad
\sum_k\frac{\pi_k}{\lambda_k}=\tfrac98,
\qquad \boxed{\Cov(M,S)=2-\tfrac{81}{32}=-\tfrac{17}{32}}.
\]
Nonzero covariance rules out unconditional independence. Classes with higher
rates tend to have higher counts but shorter waiting times; the hidden class
links the sensors after it is marginalized out.
\selfcredit{0.5 p for the expectation/covariance derivation; 0.5 p for the value and the independence conclusion.}

\solution{1e}{MAP with correlated weight regularization}{2}
The label likelihood is $\prod_{n,k}y_{nk}^{t_{nk}}$. Bayes' rule gives
\[
p(W\mid\Phi,T)\propto
\left(\prod_{n,k}y_{nk}^{t_{nk}}\right)
\exp\!\left[-\frac\alpha2\sum_kw_k^\top Rw_k\right].
\]
The evidence and Gaussian normalizers do not depend on $W$. Minimize
\[
\boxed{L(W)=-\sum_{n,k}t_{nk}\log y_{nk}
+\frac\alpha2\sum_kw_k^\top Rw_k}.
\]
Let $a_{nk}=\sum_rw_{kr}\phi_{nr}$. In index notation,
\[
\frac{\partial\log y_{nk}}{\partial w_{jr}}
=\frac{\partial}{\partial w_{jr}}
\left(a_{nk}-\log\sum_\ell e^{a_{n\ell}}\right)
=(\delta_{kj}-y_{nj})\phi_{nr}.
\]
Since each target row sums to one, the likelihood contribution is
\[
-\sum_{n,k}t_{nk}(\delta_{kj}-y_{nj})\phi_{nr}
=\sum_n(y_{nj}-t_{nj})\phi_{nr}.
\]
For the prior contribution, symmetry of $R$ is essential:
\begin{align*}
\frac{\partial}{\partial w_{jr}}\frac\alpha2
\sum_{p,q}w_{jp}R_{pq}w_{jq}
&=\frac\alpha2\left(\sum_qR_{rq}w_{jq}+\sum_pw_{jp}R_{pr}\right)\\
&=\alpha\sum_qR_{rq}w_{jq}.
\end{align*}
Therefore
\[
\boxed{\frac{\partial L}{\partial w_{jr}}
=\sum_n(y_{nj}-t_{nj})\phi_{nr}+\alpha\sum_qR_{rq}w_{jq}},
\]
\[
\boxed{\nabla_{w_j}L=\sum_n(y_{nj}-t_{nj})\phi_n+\alpha Rw_j\in\R^3}.
\]
The numerator-layout derivative is the transpose, a $1\times3$ row.
\selfcredit{0.5 p for the objective; 1 p for the component calculation including the full $R$; 0.5 p for the assembled gradient.}

\solution{1f}{Softmax ambiguity and the MAP representative}{1}
For a common shift $v$, each logit gains $v^\top\phi_n$:
\[
\frac{e^{(w_k+v)^\top\phi_n}}{\sum_je^{(w_j+v)^\top\phi_n}}
=\frac{e^{w_k^\top\phi_n}}{\sum_je^{w_j^\top\phi_n}}.
\]
The likelihood cannot distinguish these weight configurations. At a MAP
minimizer, summing the $K$ zero-gradient conditions gives
\[
0=\sum_n\left(\sum_jy_{nj}-\sum_jt_{nj}\right)\phi_n
+\alpha R\sum_jw_j=\alpha R\sum_jw_j.
\]
Since $\alpha>0$ and $R$ is invertible,
\[
\boxed{\sum_jw_j=0}.
\]
Equivalently, within any common-shift family the quadratic penalty is smallest
at $v=-K^{-1}\sum_jw_j$. The prior chooses the centered representative; it
does not remove the invariance of the prediction formula itself. Here the
positive definite penalty also makes the complete objective strongly convex
and coercive, so it has a unique finite MAP minimizer.
\selfcredit{0.5 p for shift invariance; 0.5 p for the zero-sum proof and its interpretation.}

\newpage
\section*{Exercise 2: PCA with an offset and a forbidden direction}
\solution{2a}{All affine offsets and the trace objective}{2}
Write $P=uu^\top$. Then $P^\top=P$, $P^2=P$, and
\[
x_i-\widehat x_i=(I-P)(x_i-c).
\]
Decompose $x_i-c=(x_i-\bar x)+(\bar x-c)$. The cross term disappears because
$\sum_iq_i(x_i-\bar x)=0$, so
\begin{align*}
E(c,u)
&=\sum_iq_i(x_i-\bar x)^\top(I-P)(x_i-\bar x)
+\|(I-P)(\bar x-c)\|^2\\
&=\tr[S(I-P)]+\|(I-P)(\bar x-c)\|^2.
\end{align*}
For fixed $u$, all minimizers satisfy $(I-P)(\bar x-c)=0$. Therefore
\[
\boxed{c=\bar x+\gamma u,\qquad \gamma\in\R}.
\]
These offsets describe the same affine line. In particular, $c=\bar x$ is a
valid representative but is not the only one. After minimizing over $c$,
\[
\boxed{\min_c E(c,u)=\tr[S(I-uu^\top)]=\tr(S)-u^\top Su}.
\]
Thus maximizing retained variance minimizes reconstruction error. The repetition
weights are required: they represent the full $N$-observation dataset, not an
equally weighted dataset of the $J$ distinct vectors.
\selfcredit{1 p for the decomposition and all offsets; 1 p for the trace reduction and offset interpretation.}

\solution{2b}{Unconstrained component and a new input}{1}
An orthonormal eigensystem of $S$ is
\[
u_1=\frac1{\sqrt2}\begin{pmatrix}1\\1\\0\end{pmatrix},\quad \lambda_1=7;
\qquad u_2=\frac1{\sqrt2}\begin{pmatrix}1\\-1\\0\end{pmatrix},\quad \lambda_2=3;
\qquad u_3=\begin{pmatrix}0\\0\\1\end{pmatrix},\quad \lambda_3=1.
\]
For a unit vector $u=\sum_j\beta_ju_j$, $u^\top Su=\sum_j\lambda_j\beta_j^2$
and $\sum_j\beta_j^2=1$. The largest possible retained variance is 7,
attained at $u=\pm u_1$. Choosing $u=u_1$,
\[
\boxed{\text{retained variance}=7,\qquad E_{\min}=11-7=4}.
\]
The new centered input is $r_*=x_*-\bar x=(2,1,2)^\top$. Consequently
\[
z_*=u_1^\top r_*=\frac3{\sqrt2},\qquad
\widehat x_*=\bar x+u_1z_*=\boxed{(\tfrac{11}{2},-\tfrac12,1)^\top}.
\]
Its residual is $(\tfrac12,-\tfrac12,2)^\top$, with squared norm
$\boxed{9/2}$. A single point's error need not equal the training average 4.
Changing the sign of $u_1$ changes $z_*$ but not the reconstruction.
\selfcredit{0.5 p for the direction, variance and average error; 0.5 p for the coordinate, reconstruction and pointwise error.}

\solution{2c}{The forbidden direction changes the optimum}{2}
Here $v=u_1$. Maximize $u^\top Su$ subject to both constraints with
\[
\mathcal L(u,\lambda,\eta)=u^\top Su-\lambda(u^\top u-1)-2\eta v^\top u.
\]
The stationary conditions are
\[
Su-\lambda u-\eta v=0,\qquad u^\top u=1,\qquad v^\top u=0.
\]
Multiplying the first equation by $v^\top$ and using $Sv=7v$ gives
$\eta=0$. The feasible plane is spanned by $u_2,u_3$. Their eigenvalues are
distinct, so the stationary directions and multipliers are exactly
\[
\boxed{u=\pm u_2,\ \lambda=3,\ \eta=0;
\qquad u=\pm u_3,\ \lambda=1,\ \eta=0}.
\]
Signs or multiplier values can change if a different Lagrangian sign convention
is used; the directions and their classification cannot. On the feasible unit
circle, writing $u=\beta_2u_2+\beta_3u_3$ gives variance
$3\beta_2^2+\beta_3^2\leq3$. Thus $\pm u_2$ are global maxima, while
$\pm u_3$ are global minima of retained variance.
\[
\boxed{E_{\min,\mathrm{constrained}}=11-3=8}.
\]
Using $u_2$ for the new input gives $z_*=1/\sqrt2$ and
\[
\widehat x_*=\bar x+(\tfrac12,-\tfrac12,0)^\top
=\boxed{(\tfrac92,-\tfrac52,1)^\top}.
\]
The residual is $(\tfrac32,\tfrac32,2)^\top$, so its squared norm is
$\boxed{17/2}$.
\selfcredit{1 p for the constrained stationary calculation and all directions; 0.5 p for the global maximum and average error; 0.5 p for the new-point calculation.}

\solution{2d}{PPCA likelihood and residual noise}{2}
In the eigenbasis $(u_1,u_2,u_3)$, $C$ has eigenvalues $a+\tau,\tau,\tau$.
For fixed $\mu=\bar x$, the log-likelihood is
\begin{align*}
\ell(a,\tau)
&=-\frac{3N}{2}\log(2\pi)-\frac N2\log|C|
-\frac N2\tr(C^{-1}S)\\
&=\mathrm{const}-\frac N2\left[
\log(a+\tau)+2\log\tau+\frac7{a+\tau}+\frac4\tau\right].
\end{align*}
To avoid confusing the noise variance with its standard deviation, $\tau$
itself is the variance here. Introduce $\nu=a+\tau$, so $\nu\geq\tau>0$.
The derivatives with respect to the independent coordinates are
\[
\frac{\partial\ell}{\partial\nu}
=-\frac N2\left(\frac1\nu-\frac7{\nu^2}\right),\qquad
\frac{\partial\ell}{\partial\tau}\bigg|_\nu
=-\frac N2\left(\frac2\tau-\frac4{\tau^2}\right).
\]
They vanish at $\nu=7$ and $\tau=2$, which satisfy the domain constraint.
Each bracketed scalar objective decreases up to this point and increases after
it; these are global minima of the negative log-likelihood terms. Hence
\[
\boxed{a=5,\qquad \tau=2,\qquad w=\sqrt5\,u_1}.
\]
The fitted covariance is
\[
\boxed{C=5u_1u_1^\top+2I_3
=\begin{pmatrix}9/2&5/2&0\\5/2&9/2&0\\0&0&2\end{pmatrix}}.
\]
Its eigenvalues are $(7,2,2)$, not the sample eigenvalues $(7,3,1)$.
The leading variance is fitted exactly. In the two discarded directions,
isotropic noise must use a common variance, their average $(3+1)/2=2$.
This model therefore need not reproduce the full sample covariance.
\selfcredit{0.5 p for the likelihood; 1 p for the optimizing calculation and domain check; 0.5 p for $C$ and the discrepancy.}

\solution{2e}{A different latent coordinate is not a richer marginal family}{1}
Take the positive choice of $w$ from 2d and set
\[
\boxed{w'=\frac13w=\frac{\sqrt5}{3}u_1,\qquad
\mu'=\bar x-\frac23w}.
\]
Since $\E[z']=2$ and $\Var(z')=9$, the new marginal has
\[
\E[x]=2w'+\mu'=\bar x,\qquad
\Cov(x)=9w'w'^\top+2I_3=ww^\top+2I_3=C.
\]
Both constructions are Gaussian, so equal means and covariances establish
equal marginal distributions. The negative loading choice is also valid if
$\mu'$ is adjusted to keep $\mu'+2w'=\bar x$.

Changing a scalar Gaussian latent prior cannot make the model match $S$:
its signal covariance remains rank at most one, and two marginal eigenvalues
must equal the isotropic noise variance. All three eigenvalues of $S$ are
distinct. More specifically, with the fixed variance 2,
$S-2I$ has eigenvalues $(5,1,-1)$, so it cannot be a positive semidefinite
signal covariance of any Gaussian loading. The reparameterization changes
latent coordinates, not the marginal model's expressive power.
\selfcredit{0.5 p for parameters and preserved marginal; 0.5 p for the rank/eigenvalue argument.}

\section*{Scope and source alignment}
The format follows \path{exams/homework-based-exam-1.pdf}: two extended
exercises, two hours, 18 points, derivations, and follow-up insights. The
actual second midterm's question count and topic selection are unknown.
\begin{itemize}
  \item Exercise 1 develops HW 2 Exercises 1 and 3, practice sets 3--4, and lectures 5--7. New elements are a shared rate across two different sensor distributions, a sensor missing at prediction time, marginal dependence, a full precision matrix, and the softmax common-shift ambiguity.
  \item Exercise 2 develops HW 2 Exercises 4--5, practice set 5, and lecture 10. New elements are repetition weights, all affine offsets, a forbidden direction, new numerical data, fitting the PPCA likelihood in eigenvalue coordinates, and the difference between reparameterization and increased model capacity.
\end{itemize}
The foundational operations necessarily recur, but no original homework or
practice-set problem, dataset, or complete derivation sequence is reproduced.
The questions are intentionally harder than the first midterm; the point total
does not imply equal workload.
\end{document}
